\documentclass[11pt,a4paper]{article}

% ==================================================
% Sprache und Schrift
% ==================================================

\usepackage{fontspec}
\usepackage[ngerman]{babel}

% Schrift
\setmainfont{Arial}

% ==================================================
% Seitenformat und Seitenränder
% ==================================================

\usepackage[
    a4paper,
    landscape,
    left=2cm,
    right=2cm,
    top=2cm,
    bottom=2cm
]{geometry}

% ==================================================
% Tabellen
% ==================================================

\usepackage{array}
\usepackage{longtable}
\usepackage{booktabs}
\usepackage{colortbl}
\usepackage{xcolor}

% ==================================================
% Farben
% ==================================================

\definecolor{headerblue}{RGB}{45,75,110}
\definecolor{lightblue}{RGB}{235,242,248}

% ==================================================
% Hyperlinks
% ==================================================

\usepackage[
    colorlinks=true,
    linkcolor=blue,
    urlcolor=blue,
    citecolor=blue
]{hyperref}

% ==================================================
% Eigene Befehle für das Glossar
% ==================================================

% Verlinkt einen Begriff aus dem Text mit
% dem entsprechenden Glossareintrag.

\newcommand{\gls}[1]{%
    \hyperlink{glossar:#1}{\textcolor{blue}{#1}}%
}

% ==================================================
% Dokumenteinstellungen
% ==================================================

\setlength{\parindent}{0pt}
\setlength{\parskip}{0.5em}

% ==================================================
% Dokument
% ==================================================

\begin{document}

% ==================================================
% TITEL
% ==================================================

\begin{center}

    {\LARGE\bfseries Exzerpt}

    \vspace{0.3cm}

    {\large Zusammenfassung und Reflexion eines Textes}

\end{center}

\vspace{0.8cm}

% ==================================================
% EXZERPTKOPF
% ==================================================

\section*{Exzerptkopf}

\renewcommand{\arraystretch}{1.5}

\begin{tabular}{@{}p{5cm}p{21cm}@{}}

\textbf{Thema}
&
\\[1.2cm]

\textbf{Bibliographische Angaben}
&
\\[2cm]

\textbf{Anmerkungen}
&
\\[2.5cm]

\end{tabular}

\vspace{0.5cm}

% ==================================================
% EXZERPT
% ==================================================

\section*{Exzerpt}

\small

\renewcommand{\arraystretch}{1.5}

% --------------------------------------------------
% Vier Spalten:
%
% Seite               = 2 cm
% Thema               = 4.5 cm
% Wichtige Aussage    = 9.5 cm
% Eigene Gedanken     = 9.5 cm
%
% --------------------------------------------------

\begin{longtable}{@{}p{2cm}p{4.5cm}p{9.5cm}p{9.5cm}@{}}

% ==================================================
% Tabellenkopf
% ==================================================

\toprule

\rowcolor{headerblue}

\color{white}\textbf{Seite}
&
\color{white}\textbf{Thema}
&
\color{white}\textbf{Wichtige Aussage}
&
\color{white}\textbf{Eigene Gedanken}

\\

\midrule

\endfirsthead

% ==================================================
% Tabellenkopf auf Folgeseiten
% ==================================================

\toprule

\rowcolor{headerblue}

\color{white}\textbf{Seite}
&
\color{white}\textbf{Thema}
&
\color{white}\textbf{Wichtige Aussage}
&
\color{white}\textbf{Eigene Gedanken}
\\

\midrule

\endhead

% ==================================================
% Tabellenfuß
% ==================================================

\bottomrule

\endfoot

% ==================================================
% BEISPIELZEILEN
% ==================================================

12
&
Textverständnis
&
Die \gls{Hermeneutik} beschäftigt sich mit
der Interpretation und dem Verstehen von Texten.
Dabei werden sowohl der Text selbst als auch
sein historischer und kultureller Zusammenhang
berücksichtigt.
&
Der Begriff ist insbesondere für die Interpretation
historischer Texte relevant. Interessant wäre ein
Vergleich mit modernen Methoden der
Textinterpretation.

\\[2.5cm]

13
&
Wissenschaftstheorie
&
Ein wissenschaftliches \gls{Paradigma}
bestimmt, welche Fragen innerhalb eines
bestimmten wissenschaftlichen Rahmens gestellt
werden können und welche Methoden dabei
als angemessen gelten.
&
Hier stellt sich die Frage, inwieweit sich
Paradigmen historisch verändern und welche
Folgen ein Paradigmenwechsel für die
wissenschaftliche Forschung hat.

\\[2.5cm]

14
&
Sprachliche Auseinandersetzung
&
Der Autor verwendet den Begriff
\gls{Diskurs}, um die gesellschaftliche
Auseinandersetzung mit einem bestimmten
Thema zu beschreiben.
&
Der Begriff könnte mit anderen Formen
öffentlicher Kommunikation verglichen werden.

\\[2.5cm]

15
&
Textinterpretation
&
Die \gls{Exegese} dient dazu, die Bedeutung
eines Textes systematisch zu erschließen.
Dabei werden unterschiedliche Ebenen
des Textes berücksichtigt.
&
Der Unterschied zwischen Exegese und
allgemeiner Interpretation sollte genauer
untersucht werden.

\\[2.5cm]

16
&
&
&
\\[2.5cm]

17
&
&
&
\\[2.5cm]

18
&
&
&
\\[2.5cm]

19
&
&
&
\\[2.5cm]

20
&
&
&
\\[2.5cm]

\end{longtable}

\normalsize

% ==================================================
% GLOSSAR
% ==================================================

\newpage

\section*{Glossar}

\renewcommand{\arraystretch}{1.5}

\begin{longtable}{@{}p{5cm}p{21cm}@{}}

% ==================================================
% Tabellenkopf
% ==================================================

\toprule

\textbf{Begriff}
&
\textbf{Erläuterung}

\\

\midrule

\endfirsthead

\toprule

\textbf{Begriff}
&
\textbf{Erläuterung}

\\

\midrule

\endhead

% ==================================================
% Tabellenfuß
% ==================================================

\bottomrule

\endfoot

% ==================================================
% GLOSSARBEGRIFFE
% ==================================================

\hypertarget{glossar:Hermeneutik}{}
\textbf{Hermeneutik}
&
Lehre und Theorie des Verstehens und der Interpretation
von Texten, insbesondere unter Berücksichtigung ihres
historischen und kulturellen Zusammenhangs.

\\[1cm]

\hypertarget{glossar:Paradigma}{}
\textbf{Paradigma}
&
Grundlegendes Denkmuster oder wissenschaftlicher
Bezugsrahmen, innerhalb dessen bestimmte Fragen
und Probleme betrachtet werden.

\\[1cm]

\hypertarget{glossar:Diskurs}{}
\textbf{Diskurs}
&
Geordnete Form sprachlicher Auseinandersetzung über
einen bestimmten Gegenstand oder Themenbereich.

\\[1cm]

\hypertarget{glossar:Exegese}{}
\textbf{Exegese}
&
Auslegung und Interpretation eines Textes mit dem Ziel,
dessen Bedeutung systematisch zu erschließen.

\\[1cm]

\end{longtable}

\end{document}
